\chapter{Software Architecture Design}
\label{chap:software-architecture-design}

\section{System Architecture Overview}
\label{section:architecture-overview}

\begin{figure}[htbp]
    \centering
    \includegraphics[width=0.8\textwidth]{diagrams/architecture.png}
    \caption{High-level Software Architecture Diagram}
    \label{fig:use-case-diagram}
\end{figure}

\usevar{\srsNickName} utilize pipeline architecture for analysis workflow allowing its function to be change, add, and alter its analysis.
The application lifecycle will be managed by separating the user interface, controller and pipeline control unit.
The user interface will be responsible for receiving user input and displaying results,
while the controller will manage the flow of data between the user interface and the pipeline control unit.
The pipeline control unit will orchestrate the execution of various analysis components, allowing for flexibility and modularity in the system design.

Key Software Architectural Components:

\begin{itemize}
    \item \textbf{UserInterface:} Responsible for receiving user input and displaying results, providing an intuitive interface for users to interact with the application.
    \item \textbf{Controller:} Manages the flow of data between the user interface and the pipeline control unit, ensuring that user inputs are correctly processed and results are displayed appropriately.
    \item \textbf{PipelineControlUnit:} Invokes and manage execution of pipeline nodes, allowing for dynamic configuration of analysis workflows.
    \item \textbf{ContextStorage:} Responsible for safely storing and managing state information and intermediate data (such as the TOML IR and Graph structures) throughout the analysis process, ensuring that relevant data is accessible to the Controller without interrupting the pipeline.
\end{itemize}

\section{Software Design}
\label{section:software-design}
<TIP: Describe the runtime behavior of your system.
Include a Sequence Diagram for your most important use cases.
Focus on scenarios with non-trivial interactions between components.
If your project involves a novel algorithm,
describe it here and explain why you designed it this way. />

\section{AI and Data Design}
\label{section:ai-data-design}

\subsection{Data Sources and Collection}
<TIP: Describe where your data comes from. Is it collected from users,
scraped from the web, sourced from a public dataset, or
provided by an external API?
Explain how much data you have or plan to collect,
and any challenges you anticipate in obtaining it.
Explain how data is cleaned, transformed, or prepared before use. />
\subsection{Model Design}
<TIP: Describe the AI model or external AI service you plan to use.
Explain why you chose it over alternatives.
This can connect back to your competitor or literature review in Chapter 2.
If you are training your own model, describe the architecture and
training approach. If you are using a third-party model or API,
describe how it is integrated into your system and any prompt design
or fine-tuning involved. />

\textbf{Model choice}
\\

For the classification model, we use a GraphSAGE model for architectural smell classification.
From exsisting literature, it has been shown that inductive models outperform transitive models
in code classification tasks {insert ref here}. While GATs can be modelled as an inductive algorithm,
they are prone to overfitting on graph attributes causing poor performance on unseen data. Additionally, 
when comparing GraphSAGE to GCN and GAT models, GraphSAGE is significantly faster on graphs with large edges and
nodes. It is up to 7 and 44 times faster than GAT and GCN respectively {insert ref here}.
\\

\textbf{Training methodology}
\\

We use a supervised learning technique for GNN architectural smell classification.
First, we select N projects from the Q-Scored dataset with a mininm of 100 projects for
each architectural smell. 
\\


\subsection{Evaluation}
<TIP: Describe the metrics you plan to use to evaluate your AI model. />

For the GraphSAGE evaluation, we use the following metrics
for each architectural smell:
\\
\begin{itemize}
    \item \textbf{Accuracy:} Percentage of correctly identified smells
    \item \textbf{F1:} Harmonic mean of Precision and recall
    \item \textbf{Precision:} Ratio of True positive and False positive
    \item \textbf{Recall:} Ratio of True positive and True negatives
    \item \textbf{AUC - ROC:}
\end{itemize}


For evaluating LLM explanation quality, we aim to conduct a survey with
30 participants to evaluate model explanation quality. Each survey question
is scaled from 1 - 5, with 5 being most satisfactory and 1 being least satisfactory.
The survey contains the following questions:
\\

\begin{enumerate}
  \item Is the explanation relavent to the architectural smell
  \item How accurate is the model's explanation for the smell
  \item Compared to one-shot prompting, is the model's explanation quality better.
  \item How satisfied are you with the model's explanation.
\end{enumerate}
